\chapter{KẾT LUẬN VÀ KIẾN NGHỊ}
\section{Kết luận}
Tổng hợp kết quả thực nghiệm từ Chương 3 và Chương 4. Trả lời từng câu hỏi nghiên cứu, phân biệt kết quả mô tả, kiểm định và dự báo.
\begin{itemize}[leftmargin=1.2cm]
\item Mối liên hệ giữa thời lượng dùng điện thoại và chỉ số giấc ngủ: [Điền].
\item Kết quả về độ sáng và bộ lọc ánh sáng xanh: [Điền].
\item Kết quả về caffeine và mệt mỏi: [Điền].
\item Khác biệt theo nghề nghiệp và ứng dụng: [Điền].
\item Mô hình phù hợp nhất theo tiêu chí đã chọn: [Điền].
\item Đặc trưng nổi bật theo phân tích giải thích: [Điền].
\end{itemize}
\section{Đóng góp của nghiên cứu}
Nêu đóng góp thực tế trong kiểm tra dữ liệu, phân tích thống kê, xây dựng quy trình đánh giá mô hình và diễn giải dự đoán. Chỉ khẳng định những đóng góp có bằng chứng trong báo cáo.
\section{Khuyến nghị}
Khuyến nghị phải dựa trên kết quả và tài liệu tham khảo. Nếu dữ liệu chỉ cho thấy tương quan, dùng cách diễn đạt “có liên hệ với” thay vì “gây ra”. Không đề xuất ngưỡng sử dụng điện thoại hoặc caffeine cụ thể nếu chưa có bằng chứng phù hợp.
\section{Hạn chế}
Thảo luận nguồn gốc và tính đại diện của dữ liệu, cách đo lường biến, yếu tố gây nhiễu, mất cân bằng lớp, nguy cơ rò rỉ nhãn, lựa chọn mô hình và giới hạn suy luận nhân quả. Nếu dữ liệu là dữ liệu mô phỏng/tổng hợp, phải nêu rõ.
\section{Hướng nghiên cứu tiếp theo}
Có thể mở rộng bằng dữ liệu theo thời gian, kiểm định trên bộ dữ liệu độc lập, bổ sung biến gây nhiễu, thử nghiệm mô hình khác và đánh giá độ ổn định của phương pháp giải thích.
